% N2ME_theory.tex — working note (2026-07-04)
% Bi-critical neighbourhood for the green interval JSP:
% definitions, inherited results, and the bi-objective connectivity theorem
% (what is proven, what is open). Companion to the ASOC tabu paper's
% H(sigma) framework; notation follows that paper.
\documentclass[11pt,a4paper]{article}
\usepackage[margin=2.4cm]{geometry}
\usepackage{amsmath,amssymb,amsthm}
\usepackage[T1]{fontenc}
\usepackage{microtype}

\newtheorem{theorem}{Theorem}
\newtheorem{lemma}{Lemma}
\newtheorem{proposition}{Proposition}
\newtheorem{definition}{Definition}
\newtheorem{remark}{Remark}
\newtheorem{openq}{Open question}

\newcommand{\Cmax}{C_{\max}}
\newcommand{\NPE}{\mathrm{NPE}}
\title{\textbf{A bi-critical neighbourhood for the green interval JSP}\\
{\large Working note: definitions, inherited guarantees, and the
bi-objective connectivity theorem}}
\author{H.~D\'iaz}
\date{July 2026 --- draft for discussion}

\begin{document}
\maketitle

\section{Setting and notation}

We consider the interval JSP as in the ASOC tabu paper. A processing order
$\sigma\in\Sigma$ orients the disjunctive arcs $D(\sigma)$ of the solution
graph $G(\sigma)$; the extreme graphs $G^{-}(\sigma)$ and $G^{+}(\sigma)$
weight each arc $(x,y)$ with $p_x^{-}$ and $p_x^{+}$ respectively. For
$i\in\{-,+\}$, the \emph{head} $r^i_v$ of a task $v$ is its earliest starting
time in $G^{i}(\sigma)$ (semi-active timing, computed component-wise), i.e.
\[
r^i_v=\max\bigl\{\,r^i_{JP(v)}+p^i_{JP(v)},\; r^i_{MP_\sigma(v)}+p^i_{MP_\sigma(v)}\,\bigr\},
\]
with the convention that a missing predecessor contributes $0$. Here $JP(v)$
and $MP_\sigma(v)$ denote the job and machine predecessors of $v$.

\begin{definition}[Tight arc, sustaining chain]\label{def:tight}
An arc $(x,y)$ of $G(\sigma)$ (conjunctive or disjunctive) is
\emph{$i$-tight} if $r^i_y=r^i_x+p^i_x$. A \emph{sustaining chain of $v$ in
component $i$}, written $\pi^i(v)$, is any chain
$w_1\to w_2\to\cdots\to w_r=v$ of $i$-tight arcs with $r^i_{w_1}=0$.
By construction of the heads, every task has at least one sustaining chain,
and $r^i_v$ equals the length of any such chain.
\end{definition}

Sustaining chains generalise critical paths: a critical path of
$G^{i}(\sigma)$ is precisely a sustaining chain of a task realising
$C^i_{\max}$. All makespan-criticality notions of the ASOC paper are
recovered as the special case $v\in\arg\max$.

\section{Objectives}

The primary objective is the interval makespan
$\mathbf{C}_{\max}(\sigma)=[C^-_{\max},C^+_{\max}]$ ranked by LEX2. The
secondary objective is the semi-active \emph{non-processing energy}. For a
machine $k$ let $u_1^k\prec u_2^k\prec\cdots\prec u_{n_k}^k$ be its
processing sequence under $\sigma$; the component-$i$ gap between
consecutive tasks is
$\gamma^i(u_j^k,u_{j+1}^k)=r^i_{u_{j+1}^k}-\bigl(r^i_{u_j^k}+p^i_{u_j^k}\bigr)\ge 0$,
and
\begin{equation}\label{eq:npe}
\NPE^i(\sigma)\;=\;\sum_{k\in M} P_k \sum_{j<n_k}\gamma^i\bigl(u^k_j,u^k_{j+1}\bigr)
\;=\;\sum_{k\in M} P_k\Bigl(\underbrace{r^i_{u^k_{n_k}}+p^i_{u^k_{n_k}}}_{=:~C^i(k)}
-\underbrace{r^i_{u^k_1}}_{=:~S^i(k)}\Bigr)\;-\;\text{const},
\end{equation}
where $P_k$ is the passive power of machine $k$ and the constant is
$\sum_k P_k\sum_{v\in k}p^i_v$, independent of $\sigma$. Identity
\eqref{eq:npe} (``gap form'' = ``window form'' minus a constant) is the
key structural tool below: NPE improves either by \emph{reducing a machine
completion} $C^i(k)$ or by \emph{delaying a machine start} $S^i(k)$.

\begin{remark}[Non-regularity]\label{rem:nonregular}
$\NPE$ is not a regular measure: with free timing it can be reduced by
right-shifting tasks at constant $\Cmax$. Throughout this note $\NPE$ is
evaluated on the \emph{semi-active} schedule, which makes it a function of
the order $\sigma$ alone. Consequently, no analogue of the ASOC
Proposition~2 (``only critical reversals can improve'') can hold verbatim
for free-timing NPE; the statements below are relative to semi-active
timing, and the timing dimension is delegated to a separate right-shift
operator.
\end{remark}

\section{The neighbourhood}

Let $N_2(\sigma)$ denote the published makespan neighbourhood: reversals of
boundary arcs of extreme critical blocks (arcs on critical paths of
$G^{-}$ or $G^{+}$).

\begin{definition}[Energy-critical arcs]\label{def:ecrit}
For $i\in\{-,+\}$, let $\mathcal{V}^i(\sigma)$ contain (i) every task $v$
that is the right endpoint of a positive gap
($\gamma^i(u,v)>0$ for its machine predecessor $u$), and (ii) the last task
$u^k_{n_k}$ of every machine $k$. The \emph{energy-critical arcs} are the
disjunctive arcs lying on some sustaining chain of some $v\in
\mathcal{V}^i(\sigma)$:
\[
E(\sigma)\;=\;\bigl\{(x,y)\in D(\sigma) : \exists\, i\in\{-,+\},\;
\exists\, v\in\mathcal{V}^i(\sigma),\; (x,y)\in\pi^i(v)\bigr\}.
\]
\end{definition}

\begin{definition}[Bi-critical neighbourhood]\label{def:n2me}
$N_2^{ME}(\sigma)=\{\sigma_{(a)} : a\in N_2\text{-arcs}(\sigma)\cup
E(\sigma)\}$. In \emph{lexicographic mode} the energy moves are accepted
only if the resulting makespan is not LEX2-worse; in \emph{Pareto mode}
all moves are offered to the acceptance criterion.
\end{definition}

The rationale for taking whole sustaining chains rather than only
gap-adjacent arcs comes from the connectivity analysis (Lemma~\ref{lem:C}
fails for the gap-adjacent subset).

\section{Inherited guarantees}

\begin{theorem}[Feasibility]\label{thm:feas}
The reversal of any $i$-tight disjunctive arc yields a feasible order.
In particular every neighbour in $N_2^{ME}(\sigma)$ is feasible.
\end{theorem}

\begin{proof}[Proof sketch]
The ASOC feasibility proof (Theorem~1 there) uses only tightness: if a
cycle appeared after reversing $(x,y)$, there would be a path $\pi$ from
$x$ to $y$ in $G(\sigma)$ avoiding $(x,y)$, of length
$\ell(\pi)>p^i_x$ in $G^i$; tightness $r^i_y=r^i_x+p^i_x$ then gives
$r^i_y\ge r^i_x+\ell(\pi)>r^i_y$, a contradiction. Arcs of $N_2$ are tight
because critical paths are sustaining chains; arcs of $E(\sigma)$ are tight
by Definition~\ref{def:ecrit}.
\end{proof}

\begin{theorem}[Connectivity to the makespan-optimal set — inherited]
\label{thm:inherited}
$N_2^{ME}\supseteq N_2$, hence from every $\sigma$ there is a finite
sequence of $N_2^{ME}$ moves reaching some makespan-optimal order
(ASOC Theorem~2).
\end{theorem}

This is all that the superset argument yields. It says nothing about
reaching the \emph{lexicographically} optimal order among the
makespan-optimal ones, which is the object of the next section.

\section{Bi-objective connectivity}

Fix a lexicographic optimum $\sigma^{\star}$ (minimal
$\mathbf{C}_{\max}$ under LEX2; among those, minimal
$[\NPE^-,\NPE^+]$ under LEX2). Following the inversion-count scheme of
ASOC Theorem~2, connectivity to $\sigma^{\star}$ reduces to a
\emph{key lemma}: \textbf{every non-optimal $\lambda$ owns a disagreeing
arc — $(x,y)\in D(\lambda)$, $(y,x)\in D(\sigma^{\star})$ — inside
$N_2^{ME}(\lambda)$}. Reversing it decreases the inversion count
$M(\lambda,\sigma^{\star})$, and induction finishes the proof. Three cases:

\medskip\noindent\textbf{Case A (makespan deficit).}
$\mathbf{C}_{\max}(\lambda)$ LEX2-worse than
$\mathbf{C}_{\max}(\sigma^{\star})$. Covered verbatim by the ASOC key
lemma: some extreme critical path of $\lambda$ carries a disagreeing arc,
which belongs to the $N_2$ component. \hfill$\checkmark$

\begin{lemma}[Case B: completion-side energy deficit]\label{lem:C}
Suppose $\mathbf{C}_{\max}(\lambda)=\mathbf{C}_{\max}(\sigma^{\star})$ and
there exist $i$ and a machine $k$ with $C^i(k)(\lambda)>C^i(k)(\sigma^{\star})$,
where the machine sequences of $\lambda$ and $\sigma^\star$ on $k$ end in
the same task\footnote{The general statement compares
$\max_{v\in k}(r^i_v+p^i_v)$, which is order-independent as a set of
candidate tasks; the footnoted condition is dispensable at the price of a
longer case analysis.}. Then some sustaining chain
$\pi^i\bigl(u^k_{n_k}\bigr)$ in $\lambda$ contains a disagreeing arc, and
that arc lies in $E(\lambda)$.
\end{lemma}

\begin{proof}
Let $\pi$ be a sustaining chain of $v=u^k_{n_k}$ in $\lambda$, of length
$C^i(k)(\lambda)-p^i_v$ in $G^i(\lambda)$. If every disjunctive arc of
$\pi$ agreed with $\sigma^{\star}$, then $\pi$ would be a path of
$G^i(\sigma^{\star})$, whence
$r^i_v(\sigma^{\star})\ge \ell(\pi)=r^i_v(\lambda)$ and
$C^i(k)(\sigma^{\star})\ge C^i(k)(\lambda)$, contradicting the hypothesis.
The disagreeing arc is on a sustaining chain of a machine-last task, hence
in $E(\lambda)$ by Definition~\ref{def:ecrit}(ii).
\end{proof}

\medskip\noindent\textbf{Case C (start-side energy deficit) — OPEN.}
The remaining possibility, by identity \eqref{eq:npe}, is
$C^i(k)(\lambda)=C^i(k)(\sigma^{\star})$ for all $k,i$ while some machine
start satisfies $S^i(k)(\lambda)<S^i(k)(\sigma^{\star})$: $\lambda$ switches
machine $k$ on \emph{too early}. The chain argument does not transfer: a
sustaining-chain of the first task of $k$ in $\lambda$ certifies a
\emph{lower bound} that is already met, whereas the deficit asks why
$\sigma^{\star}$ manages to start \emph{later}. The witness structure lives
in $\sigma^{\star}$, not in $\lambda$, so it cannot by itself point to an
arc of $E(\lambda)$. Two candidate resolutions, in decreasing order of
elegance:
\begin{enumerate}
\item \textbf{Enlarge $E(\sigma)$ with delay moves.} Add, for every
machine-first task $f$ preceding a positive gap, the reversal of
$(f,MS_\sigma(f))$ — a swap that sends $f$ deeper into the sequence and can
only delay the machine start. Conjecture: with these moves Case~C admits an
exchange argument on the first disagreeing position of machine $k$'s
sequence. \emph{Status: unproven.}
\item \textbf{Weaken the theorem.} State connectivity to the set
$\{\sigma:\mathbf{C}_{\max}\text{ optimal},\; C^i(k)\text{ minimal}\}$
(Cases A+B, fully proven) and treat the start-side as heuristic
intensification, validated empirically against exact CP-SAT fronts on
small instances. This mirrors the pragmatic scope of the connectivity
results in the group's energy-neighbourhood paper for the fuzzy flexible
JSP, whose Theorem~1 likewise proves reachability with respect to the
structures its moves control.
\end{enumerate}

\begin{openq}
Does Definition~\ref{def:ecrit} augmented with the delay moves of
resolution (1) satisfy the Case-C key lemma? Equivalently: can a
makespan- and completion-optimal order with a strictly early machine start
always be brought closer to $\sigma^{\star}$ by a single tight-arc reversal
or first-task swap?
\end{openq}

\section{Remarks for the paper}

\begin{remark}[Why not only gap-adjacent arcs]
If $E(\sigma)$ contained only the arcs $(u,v)$ delimiting a positive gap,
Lemma~\ref{lem:C} would fail: the disagreeing arc found in its proof lies
anywhere along the sustaining chain, typically far upstream of the gap.
This is the formal reason the neighbourhood must look at sustaining chains
— the energy analogue of looking at whole critical paths rather than the
last arc before $\Cmax$.
\end{remark}

\begin{remark}[Pareto mode]
For dominance-based search the relevant notion is connectivity to
\emph{each} Pareto-optimal order. Cases A and B apply unchanged to any
non-dominated target $\sigma^{\star}$ (a dominated $\lambda$ has a deficit
in some objective and component); Case C is open exactly as above. The
lexicographic theorem is the extreme-point instance of this statement.
\end{remark}

\begin{remark}[Relation to the N2Plus result]
Under pure single-objective LEX2, the worst-case sub-neighbourhood
(G$^{+}$ only) suffices — the unpublished N2Plus experiment. As soon as the
ranking involves a second criterion (lexicographic tie-break or Pareto
dominance), the ranking-agnostic form of the connectivity argument is
required again, over both extreme graphs and, per this note, over both
objective structures. The bi-objective setting is thus precisely where the
full bi-critical machinery earns its keep.
\end{remark}

\end{document}
